\section{Introduction}
A prediction interval on remaining useful life (RUL) becomes an engineering commitment when a maintenance rule consumes it. A planner who triggers work whenever a one-sided lower bound falls to or below a lead time of twenty cycles trusts that bound in two directions: it must not be so low that equipment far from needing attention is pulled early, and not so high that equipment that needs attention stays in service. Interval quality is usually summarised by coverage and width on clean inputs. A coverage figure does not say which of the two errors a failure produces, or how often.

Three adjacent strands of work exist. Interval and uncertainty-quantification methods for RUL, including MC Dropout, deep ensembles, heteroscedastic networks and conformal prediction, have been developed and benchmarked on C-MAPSS and N-CMAPSS.\cite{peng2019,basora2025,zhan2024,piao2025,mao2024,xu2021,xu2026} A second strand studies RUL prediction under sensor malfunction or failure, mainly reporting point-prediction accuracy.\cite{xu2023,tang2026,hachem2024} A third converts probabilistic or interval RUL into maintenance decisions and costs.\cite{kim2020,lee2022,zhu2025} A bounded literature search for this revision (Consensus, October 2026; about twenty papers, mostly abstract-level) found no study reporting interval coverage under structured sensor drift or bias together with maintenance-trigger error rates. That describes the search only; no priority is claimed.

This paper examines the intersection in a controlled stress test. Intervals come from five mechanisms on two backbones (LSTM and Transformer) for the four C-MAPSS sub-datasets, under an engine-level protocol with disjoint fitting, validation and calibration engines. At the final window of each test engine we perturb the sensor channels (additive noise, bias, drift or gain error), apply a one-sided 95\% lower-bound rule at lead times of 10, 20 and 30 cycles, and count two errors separately: engines that needed maintenance but were not triggered, and engines triggered although far from needing it. Three questions are asked. (Q1) In which direction and by how much do structured degradations break coverage, and does the share of out-of-range inputs predict the damage? (Q2) Which trigger errors result, how uncertain are they at the engine level, and does the mean or the scale output of the model drive them? (Q3) Does the cost ranking of mechanisms survive a change of cost ratio?

Three contributions are claimed, each bounded by the conditions stated in later sections. First, evidence that the side on which the lower bound errs follows the sign pattern of the sensor shift---a coherent positive shift gives premature triggers, a coherent negative shift gives unrecognised needs---and that coverage and the share of out-of-range inputs do not reveal that direction. Second, an engine-level accounting of both trigger errors with resampling intervals and a same-condition mean/scale output swap. Third, a transparent protocol record that keeps the negative results: no mechanism is recommendable on cost, the scale output has no order-independent sign effect under additive noise, and controls that change the waveform of the shift (ramp, constant offset, reversed or shuffled ramp) show sensitivity to the waveform without identifying temporal order as its cause. The study is empirical: it proposes no architecture, drift detector or maintenance policy, uses simulated data and one terminal decision per engine, and is limited accordingly. Repairs to an earlier version of the pipeline (Section~\ref{sec:framework}) are reported for transparency, not as a contribution.
